% GENERATED FILE. Edit canonical lessons, then run npm run generate:books.
% canonical-source-hash: fnv1a64:400178c8c7843e81
% canonical-lessons: JA-C270-hannin, JA-C270-jiken, JA-C270-tsuuhou, JA-C270-sono, JA-C270-sawaru

\chapter{Who Did It?}
\label{ch:ja-who-did-it}
\hlchaptermodality{270}

\begin{chapteropening}
\textbf{By the end of this chapter:} \emph{I can say what the police were told and what not to touch.}

Complete the last lesson of chapter 270: I can say what the police were told and what not to touch.
\end{chapteropening}

\section[\texorpdfstring{\ja{はんにん} (hannin)}{はんにん (hannin)}]{\ja{はんにん} (hannin) — a culprit; the criminal}

\label{lesson:JA-C270-hannin}



\begin{quote}
Before the new one: say the Japanese for to bark, then the Japanese for to make a noise.
\end{quote}

\subsection*{You'll want to know: \ja{はんにん}}
\textbf{\ja{はんにん}} — \emph{hannin} — \textquotedblleft{}a culprit; the criminal\textquotedblright{}.

\begin{quote}
\textbf{\ja{はんにんが} \ja{にげた} \ja{とき、}\ja{月}\ja{が} \ja{ひかっていました。}} — \emph{hannin ga nigeta toki, tsuki ga hikatte imashita} — \textquotedblleft{}When the culprit got away, the moon was shining.\textquotedblright{}
\end{quote}

\begin{cousinweb}
\textbf{\ja{はんにん}} (\emph{hannin}) is Sino-Japanese (kango): two Chinese characters read with their on'yomi, readings borrowed from Chinese. They are \emph{han}, \textquotedblleft{}crime\textquotedblright{} (Middle Chinese \emph{bjomX}) and \emph{nin}, \textquotedblleft{}person\textquotedblright{} (\emph{nyin}) (KANJIDIC2; Middle Chinese in Baxter's transcription).
\end{cousinweb}

\begin{grammarlens}[title={What you've built}]
One more word for telling what happened.
\end{grammarlens}

\subsection*{Your turn}
\begin{itemize}
\raggedright
  \item \emph{Say it:} \emph{hannin}
  \item \emph{Say it:} \emph{hannin}, once more
  \item \emph{Say it:} \emph{hannin ga nigeta toki, tsuki ga hikatte imashita} — When the culprit got away, the moon was shining.
  \item \emph{Recall:} say the Japanese for to bark, then the Japanese for to make a noise, then say \emph{hannin} again
\end{itemize}

\begin{tcolorbox}[breakable,colback=teal!4,colframe=teal!35!black,title={Before you move on}]
What does \ja{はんにん} mean? (\textquotedblleft{}A culprit; the criminal\textquotedblright{}.) Say it once more.
\end{tcolorbox}
\section[\texorpdfstring{\ja{じけん} (jiken)}{じけん (jiken)}]{\ja{じけん} (jiken) — an incident; a case}

\label{lesson:JA-C270-jiken}



\begin{quote}
Before the new one: say the Japanese for to make a noise, then the Japanese for a culprit.
\end{quote}

\subsection*{You'll want to know: \ja{じけん}}
\textbf{\ja{じけん}} — \emph{jiken} — \textquotedblleft{}an incident; a case\textquotedblright{}.

\begin{quote}
\textbf{\ja{この} \ja{まちで} \ja{じけんが} \ja{ありました。}} — \emph{kono machi de jiken ga arimashita} — \textquotedblleft{}There was an incident in this town.\textquotedblright{}
\end{quote}

\begin{cousinweb}
\textbf{\ja{じけん}} (\emph{jiken}) is Sino-Japanese (kango): two Chinese characters read with their on'yomi, readings borrowed from Chinese. They are \emph{ji}, \textquotedblleft{}matter\textquotedblright{} (Middle Chinese \emph{dzriH}) and \emph{ken}, \textquotedblleft{}affair\textquotedblright{} (\emph{gjenX}) (KANJIDIC2; Middle Chinese in Baxter's transcription).
\end{cousinweb}

\begin{grammarlens}[title={What you've built}]
One more word for telling what happened.
\end{grammarlens}

\subsection*{Your turn}
\begin{itemize}
\raggedright
  \item \emph{Say it:} \emph{jiken}
  \item \emph{Say it:} \emph{jiken}, once more
  \item \emph{Say it:} \emph{kono machi de jiken ga arimashita} — There was an incident in this town.
  \item \emph{Recall:} say the Japanese for to make a noise, then the Japanese for a culprit, then say \emph{jiken} again
\end{itemize}

\begin{tcolorbox}[breakable,colback=teal!4,colframe=teal!35!black,title={Before you move on}]
What does \ja{じけん} mean? (\textquotedblleft{}An incident; a case\textquotedblright{}.) Say it once more.
\end{tcolorbox}
\section[\texorpdfstring{\ja{つうほう} (tsūhō)}{つうほう (tsūhō)}]{\ja{つうほう} (tsūhō) — a report to the police; a tip-off}

\label{lesson:JA-C270-tsuuhou}



\begin{quote}
Before the new one: say the Japanese for a culprit, then the Japanese for an incident.
\end{quote}

\subsection*{You'll want to know: \ja{つうほう}}
\textbf{\ja{つうほう}} — \emph{tsūhō} — \textquotedblleft{}a report to the police; a tip-off\textquotedblright{}.

\begin{quote}
\textbf{\ja{ははが} \ja{けいさつに} \ja{つうほうしました。}} — \emph{haha ga keisatsu ni tsūhō shimashita} — \textquotedblleft{}My mother called the police about it.\textquotedblright{}
\end{quote}

\begin{cousinweb}
\textbf{\ja{つうほう}} (\emph{tsūhō}) is Sino-Japanese (kango): two Chinese characters read with their on'yomi, readings borrowed from Chinese. They are \emph{tsū}, \textquotedblleft{}pass through\textquotedblright{} (Middle Chinese \emph{thuwng}) and \emph{hō}, \textquotedblleft{}report\textquotedblright{} (\emph{pawH}) (KANJIDIC2; Middle Chinese in Baxter's transcription).
\end{cousinweb}

\begin{grammarlens}[title={What you've built}]
One more word for telling what happened.
\end{grammarlens}

\subsection*{Your turn}
\begin{itemize}
\raggedright
  \item \emph{Say it:} \emph{tsūhō}
  \item \emph{Say it:} \emph{tsūhō}, once more
  \item \emph{Say it:} \emph{haha ga keisatsu ni tsūhō shimashita} — My mother called the police about it.
  \item \emph{Recall:} say the Japanese for a culprit, then the Japanese for an incident, then say \emph{tsūhō} again
\end{itemize}

\begin{tcolorbox}[breakable,colback=teal!4,colframe=teal!35!black,title={Before you move on}]
What does \ja{つうほう} mean? (\textquotedblleft{}A report to the police; a tip-off\textquotedblright{}.) Say it once more.
\end{tcolorbox}
\section[\texorpdfstring{\ja{その} (sono)}{その (sono)}]{\ja{その} (sono) — that (one); the one mentioned}

\label{lesson:JA-C270-sono}



\begin{quote}
Before the new one: say the Japanese for an incident, then the Japanese for a report to the police.
\end{quote}

\subsection*{You'll want to know: \ja{その}}
\textbf{\ja{その}} — \emph{sono} — \textquotedblleft{}that (one); the one mentioned\textquotedblright{}.

\begin{quote}
\textbf{\ja{その} \ja{おとこを} \ja{みましたか。}} — \emph{sono otoko o mimashita ka} — \textquotedblleft{}Did you see the man?\textquotedblright{}
\end{quote}

\begin{cousinweb}
\textbf{\ja{その}} (\emph{sono}) is a native Japanese word (wago), inherited rather than borrowed from Chinese, and it is usually written in kana (JMdict).
\end{cousinweb}

\begin{grammarlens}[title={What you've built}]
One more word for telling what happened.
\end{grammarlens}

\subsection*{Your turn}
\begin{itemize}
\raggedright
  \item \emph{Say it:} \emph{sono}
  \item \emph{Say it:} \emph{sono}, once more
  \item \emph{Say it:} \emph{sono otoko o mimashita ka} — Did you see the man?
  \item \emph{Recall:} say the Japanese for an incident, then the Japanese for a report to the police, then say \emph{sono} again
\end{itemize}

\begin{tcolorbox}[breakable,colback=teal!4,colframe=teal!35!black,title={Before you move on}]
What does \ja{その} mean? (\textquotedblleft{}That (one); the one mentioned\textquotedblright{}.) Say it once more.
\end{tcolorbox}
\section[\texorpdfstring{\ja{さわる} (sawaru)}{さわる (sawaru)}]{\ja{さわる} (sawaru) — to touch}

\label{lesson:JA-C270-sawaru}



\begin{quote}
Before the new one: say the Japanese for a report to the police, then the Japanese for that (one).
\end{quote}

\subsection*{You'll want to know: \ja{さわる}}
\textbf{\ja{さわる}} — \emph{sawaru} — \textquotedblleft{}to touch\textquotedblright{}.

\begin{quote}
\textbf{\ja{なにも} \ja{さわらないで} \ja{ください。}} — \emph{nanimo sawaranaide kudasai} — \textquotedblleft{}Please do not touch anything.\textquotedblright{}
\end{quote}

Polite \textbf{\ja{さわります}}, \emph{sawarimasu}; a godan verb. Past (polite): \textbf{\ja{おとこが} \ja{かばんに} \ja{さわりました。}}, \emph{otoko ga kaban ni sawarimashita}.

\begin{cousinweb}
\textbf{\ja{さわる}} (\emph{sawaru}) is a native Japanese word (wago), inherited rather than borrowed from Chinese; its kanji was attached later as a kun'yomi (JMdict).
\end{cousinweb}

\begin{grammarlens}[title={What you've built}]
One more word for telling what happened.
\end{grammarlens}

\subsection*{Your turn}
\begin{itemize}
\raggedright
  \item \emph{Say it:} \emph{sawaru}
  \item \emph{Say it:} \emph{sawaru}, once more
  \item \emph{Say it:} \emph{nanimo sawaranaide kudasai} — Please do not touch anything.
  \item \emph{Recall:} say the Japanese for a report to the police, then the Japanese for that (one), then say \emph{sawaru} again
\end{itemize}

\begin{tcolorbox}[breakable,colback=teal!4,colframe=teal!35!black,title={Before you move on}]
What does \ja{さわる} mean? (\textquotedblleft{}To touch\textquotedblright{}.) Say it once more.
\end{tcolorbox}
